前面写出了 $p_\theta(x,z)$ 和 $q_\phi(z\mid x)$，这里把它们之间的关系再展开一下。
假设 $z$ 是未观测的潜变量，$x$ 是已经观测到的数据。
先验 $p(z)$ 表示观测 $x$ 之前对 $z$ 的分布设定，
似然 $p_\theta(x\mid z)$ 表示给定 $z$ 后如何产生 $x$。
将两者相乘，得到联合分布；将 $z$ 积分掉，则得到观测数据的边缘分布：
\[
p_\theta(x,z)=p(z)p_\theta(x\mid z),\qquad
p_\theta(x)=\int p(z)p_\theta(x\mid z)\,\mathrm dz.
\]
于是，在 $p_\theta(x)>0$ 时，贝叶斯公式给出
\[
p_\theta(z\mid x)=\frac{p_\theta(x\mid z)p(z)}{p_\theta(x)}.
\]
左边就是后验。它回答的是“已经看见这个 $x$，哪些 $z$ 更有可能产生它”。
分母对固定的 $x$ 是一个常数，但求出它仍可能需要高维积分，
所以写出了贝叶斯公式，不等于已经能够计算后验。

看一个可以算完的例子。设 $Z\sim\mathcal N(0,\tau^2)$，
给定 $Z=z$，观测 $X=z+\eta$，其中 $\eta\sim\mathcal N(0,\sigma^2)$。
观测到 $X=x$ 后，后验密度中与 $z$ 有关的指数是
\[
-\frac12\left(\frac{z^2}{\tau^2}+\frac{(x-z)^2}{\sigma^2}\right).
\]
把平方项配起来，可得
\[
Z\mid X=x\sim\mathcal N\left(
\frac{\tau^2}{\tau^2+\sigma^2}x,
\frac{\tau^2\sigma^2}{\tau^2+\sigma^2}\right).
\]
观测噪声小，后验均值就更靠近 $x$；观测噪声大，就更多地依赖先验。
但后验仍然有方差，不能把“后验均值”当成“唯一可能的原始值”。
后面去噪模型遇到的也是这个问题：同一个带噪观测可能来自不同干净样本。

另外，“似然”与“概率密度”不要求是两个不同公式。
固定参数时，$p_\theta(x)$ 是 $x$ 的密度；固定已经观测的数据 $x$，
将同一个式子看成参数 $\theta$ 的函数，就是似然。
极大似然估计选择让观测数据更可能出现的参数。
如果其中的积分难算，前面的 ELBO 就提供了可优化的下界；
$q_\phi(z\mid x)$ 是为此引入的近似后验，不是直接把未知真后验当作已知。
这些区别在把整条扩散轨迹当成潜变量时仍然适用。\cite{theorybook}
